%%%%%%%%%%%%%%%%%%%%%%%%%%%%%%%%%%%%%%%%%
% "ModernCV" CV and Cover Letter
% LaTeX Template
% Version 1.1 (9/12/12)
%
% This template has been downloaded from:
% http://www.LaTeXTemplates.com
%
% Original author:
% Xavier Danaux (xdanaux@gmail.com)
%
% License:
% CC BY-NC-SA 3.0 (http://creativecommons.org/licenses/by-nc-sa/3.0/)
%
% Important note:
% This template requires the moderncv.cls and .sty files to be in the same 
% directory as this .tex file. These files provide the resume style and themes 
% used for structuring the document.
%
%%%%%%%%%%%%%%%%%%%%%%%%%%%%%%%%%%%%%%%%%

%----------------------------------------------------------------------------------------
%	PACKAGES AND OTHER DOCUMENT CONFIGURATIONS
%----------------------------------------------------------------------------------------

\documentclass[11pt,a4paper,sans]{moderncv} % Font sizes: 10, 11, or 12; paper sizes: a4paper, letterpaper, a5paper, legalpaper, executivepaper or landscape; font families: sans or roman
\usepackage{standalone}
\moderncvstyle{classic} % CV theme - options include: 'casual' (default), 'classic', 'oldstyle' and 'banking'
\moderncvcolor{blue} % CV color - options include: 'blue' (default), 'orange', 'green', 'red', 'purple', 'grey' and 'black'


\usepackage[scale=0.85]{geometry} % Reduce document margins
%\setlength{\hintscolumnwidth}{3cm} % Uncomment to change the width of the dates column
%\setlength{\makecvtitlenamewidth}{10cm} % For the 'classic' style, uncomment to adjust the width of the space allocated to your name

%\usepackage[utf8]{inputenc}

%\usepackage{booktabs}
\usepackage{fontawesome}
\usepackage{marvosym} % For cool symbols.
%\usepackage{hyperref}



%----------------------------------------------------------------------------------------
%	NAME AND CONTACT INFORMATION SECTION
%----------------------------------------------------------------------------------------

\firstname{} % First name is included in the family-name field below
\familyname{Dr Andrew Mummery} % Your last name

% All information in this block is optional, comment out any lines you don't need
\title{\Letter\ \href{mailto:amummery@ias.edu}{amummery@ias.edu}}
\address{School of Natural Sciences}{Institute for Advanced Study, Princeton, NJ}
% \mobile{(+44) 7736346815}


%\fax{(000) 111 1113}
 
%\social{github}{stefano-bragaglia}



%\homepage{www.xyz.com/}{My Webpage}

% social link \faGithub, \faSkype, \faLinkedin,\faStackExchange, and \faStackOverflow
\extrainfo{
    \faGithub\href{https://github.com/andymummeryastro}{ GitHub} \quad
    \href{https://andymummeryastro.github.io}{Website}
    }

% \email{andrew.mummery@physics.ox.ac.uk} 


%\social[linkedin][www.linkedin.com]{name}
% The first argument is %the url for the clickable link, the second argument is the url displayed in the %template - this allows special characters to be displayed such as the tilde in this %example

% \photo[70pt][0.3pt]{picture} % The first bracket is the picture height, the second is %the thickness of the frame around the picture (0pt for no frame)
%\quote{Not Attention, Patience is all we need.}

%----------------------------------------------------------------------------------------

\newcommand{\cvdoublecolumn}[2]{%
  \cvitem[.75em]{}{%
    \begin{minipage}[t]{\listdoubleitemcolumnwidth}#1\end{minipage}%
    \hfill%
    \begin{minipage}[t]{\listdoubleitemcolumnwidth}#2\end{minipage}%
    }%
}



\usepackage{multibbl}
\newcommand\Colorhref[3][orange]{\href{#2}{\small\color{#1}#3}}


% \newcommand{\cvreference}[7]{%
%     \textbf{#1}\newline% Name
%     \ifthenelse{\equal{#2}{}}{}{\addresssymbol~#2\newline}%
%     \ifthenelse{\equal{#3}{}}{}{#3\newline}%
%     \ifthenelse{\equal{#4}{}}{}{#4\newline}%
%     \ifthenelse{\equal{#5}{}}{}{#5\newline}%
%     \ifthenelse{\equal{#6}{}}{}{\emailsymbol~\texttt{#6}\newline}%
%     \ifthenelse{\equal{#7}{}}{}{\phonesymbol~#7}}

\begin{document}
\makecvtitle % Print the CV title





%----------------------------------------------------------------------------------------
%	EDUCATION SECTION
%----------------------------------------------------------------------------------------
\section{Employment}
\cventry{2025--present}{John N. Bahcall Fellow and Member}{School of Natural Sciences}{Institute for Advanced Study}{}{Independent research on black hole accretion, tidal disruption events, and strong-gravity astrophysics.}{}
\cventry{2022--2025}{Leverhulme-Peierls Fellow}{Rudolf Peierls Centre for Theoretical Physics}{University of Oxford}{}{Independent research fellowship in black hole accretion physics.}{}
\cventry{2022--2025}{Junior Research Fellow}{Merton College}{University of Oxford}{}{Non-stipendiary junior research fellowship.}{}{}

\section{Education}

\cventry{2018--2022}{DPhil, Theoretical Physics}{University of Oxford}{Wadham College}{}
{Thesis: \textit{Illuminating tidal disruption events with a time-dependent theory of relativistic accretion discs} 
\newline Supervisor: Prof. Steven Balbus 
\newline Thesis awarded best-thesis prizes by the International Astronomical Union and the Royal Astronomical Society, and an ``exceptional'' letter of commendation by the University of Oxford.}

\cventry{2014--2018}{Master of Physics}{Department of Physics}{University of Oxford}{Wadham College}{First-class honours; highest degree mark in the year group. The final-year Master's project was subsequently published in the Monthly Notices of the Royal Astronomical Society.}






%----------------------------------------------------------------------------------------
%	Fellowships \& Awards
%----------------------------------------------------------------------------------------

\section{Fellowships \& Awards}
\cventry{2025--2030}{John N. Bahcall Fellowship}{Institute for Advanced Study}{}{}{}
\cventry{2022--2025}{Leverhulme-Peierls Fellowship}{University of Oxford}{}{}{}
\cventry{2022}{International Astronomical Union best thesis prize}{Division D: Fundamental Physics and High Energy Astrophysical Phenomena}{}{}{}
\cventry{2022}{Royal Astronomical Society Michael Penston thesis prize}{}{}{}{}
\cventry{2022}{Non-stipendiary junior research fellowship}{}{Merton College, Oxford University}{}{}
\cventry{2022}{Letter of commendation for an ``exceptional'' DPhil thesis}{}{Oxford University}{}{}
\cventry{2020}{Visiting research scholarship}{}{Princeton University}{}{}
\cventry{2019}{Keeley graduate scholarship}{}{Wadham College, Oxford University}{}{}
\cventry{2018}{Scott prize}{}{Oxford University}{}{Awarded for the top Physics degree mark in the class of 2018}
\cventry{2018}{Sukumar prize}{}{Wadham College, Oxford University}{}{Awarded for the top Physics degree mark in the class of 2018}

\section{Students}
\cvitem{PhD}{\textbf{Muryel Guolo} (Johns Hopkins $\rightarrow$ IAS). Papers led with \textbf{Mummery}:
\newline Guolo, \textbf{Mummery}, van Velzen et al.\ (2026), The wandering supermassive black hole powering the off-nuclear TDE AT2024tvd. ApJL 1004, L25. \href{https://arxiv.org/abs/2602.12272}{\textcolor{blue}{arXiv}}.
\newline Guolo, \textbf{Mummery}, van Velzen et al.\ (2025), Compact accretion disks in the aftermath of tidal disruption events: parameter inference from joint X-ray spectra and UV/optical photometry. Submitted to ApJ. \href{https://arxiv.org/abs/2510.26774}{\textcolor{blue}{arXiv}}.
\newline Guolo, \textbf{Mummery}, Ingram, Nicholl, Gezari \& Nathan (2025), A time-dependent solution for GSN 069 disk evolution and the nature of long-lived tidal disruption events. ApJ 992, 114. \href{https://arxiv.org/abs/2504.20148}{\textcolor{blue}{arXiv}}.
\newline Guolo, \textbf{Mummery}, Wevers, Nicholl, Gezari, Ingram \& Pasham (2025), The properties of the GSN 069 accretion disk from a joint X-ray and UV spectral analysis: stress-testing quasi-periodic eruption models. ApJ 985, 146. \href{https://arxiv.org/abs/2501.03333}{\textcolor{blue}{arXiv}}.
\newline Guolo \& \textbf{Mummery} (2025), The size of accretion disks from self-consistent X-ray spectra and UV/optical/NIR photometry fitting: ASASSN-14li and HLX-1 --- diskSED. ApJ 978, 167. \href{https://arxiv.org/abs/2408.17296}{\textcolor{blue}{arXiv}}.}
\cvitem{PhD}{\textbf{Jake Rule} (University of Oxford; on the job market). Selected papers:
\newline Rule, \textbf{Mummery}, Balbus, Stone \& Zhang (2025), The plunging region of a thin accretion disc around a Schwarzschild black hole. MNRAS 542, 377. \href{https://arxiv.org/abs/2505.13701}{\textcolor{blue}{arXiv}}.
\newline Rule, \textbf{Mummery}, Balbus, Stone \& Zhang (2026), The plunging region of thin accretion discs across the black hole spin range. Accepted in MNRAS. \href{https://arxiv.org/abs/2608.06278}{\textcolor{blue}{arXiv}}.}
\cvitem{PhD}{\textbf{Fangyuan (Jason) Yu} (Princeton University). Papers led with \textbf{Mummery}:
\newline Yu, \textbf{Mummery} \& Guolo (2026), HLX--1 as a repeating partial tidal disruption event around an intermediate-mass black hole. Accepted in ApJ. \href{https://arxiv.org/abs/2609.38694}{\textcolor{blue}{arXiv}}.}
\cvitem{Summer}{\textbf{Wenkang Xin}. Xin \& \textbf{Mummery} (2026), The relativistic tidal tensor: general solutions for stationary axisymmetric spacetimes and the Hills mass of naked singularities. Physical Review D 113, 044036. \href{https://arxiv.org/abs/2511.21499}{\textcolor{blue}{arXiv}}. Now a student with John Magorrian at the University of Oxford.}

\section{Observing programmes}
\cventry{2025--}{Hubble Space Telescope GO 18081 (Principal Investigator)}{}{}{}{\textit{Using 35-year-old tidal disruption events to test the fundamentals of accretion disk theory}. Award: \$60{,}000.}
\cventry{Cycle 5}{JWST GO 11908 (co-PI)}{}{}{}{\textit{Probing the $M$--$\sigma$ Relation in the Intermediate-Mass Black Hole Regime with Off-Nuclear TDEs}. Award: \$175{,}000.}
\cventry{Cycle 5}{JWST GO 12397 (co-PI)}{}{}{}{\textit{Where Are the Stars? A Search for the Stellar Counterparts of Two Off-Nuclear TDE Candidates}. Award: \$110{,}000.}
\cventry{Cycle 8}{NICER C8 9062 (co-PI)}{}{}{}{\textit{A Time-Dependent View of XRB Outbursts: Spectroscopy Beyond Steady-State Models}. Award: \$60{,}000.}


%----------------------------------------------------------------------------------------
%	Teaching SECTION
%----------------------------------------------------------------------------------------

\section{Teaching and Supervising}
\cventry{2024--2025}{Lecturer: Magnetohydrodynamics}{}{University of Oxford}{}{
% {Lecturer for the MHD course in the Masters in Mathematical and Theoretical Physics \newline 
I wrote and gave 10 one-hour lectures covering the fundamentals of MHD through to topics of modern research. This course is attended by Masters and PhD students in the Mathematics and Physics departments. I then wrote, set and marked an exam paper for the Masters students. Lecture notes available \href{https://github.com/andymummeryastro/MHD_notes}{\textcolor{blue}{here}}. }{}
\cventry{2018-2022}{Problem classes -- (College Tutor)}{}{University of Oxford (various colleges)}
{\newline -- General Relativity (3rd year undergraduate)
\newline -- Special Relativity (3rd year undergraduate)
\newline -- Chaos, Random Processes and Predictability (2nd year undergraduate)
\newline -- Fluid Dynamics (3rd year undergraduate)}{Led small group problem classes (typically 2-5 students), marked and reviewed mock exam papers, and led final exam revision sessions for larger classes (typically 8-10 students).}
\cventry{2021--2024}{Undergraduate interviews}{}{University College, University of Oxford}{}{Interviewed roughly 50 prospective undergraduates each year; wrote physics and mathematics questions and served on the College shortlisting and selection panel.}


\section{Peer review}
\cventry{2019--present}{Journal refereeing}{}{}{}{Peer reviewer for the Monthly Notices of the Royal Astronomical Society, The Astrophysical Journal, The Astrophysical Journal Letters, The Open Journal of Astrophysics, Physical Review D, Astronomy \& Astrophysics, and Astroparticle Physics.}{}{}{}

\cventry{2024}{PhD examining}{}{}{}{External PhD examiner for the thesis defence of Jesse Voss at the University of Radboud.}{}{}

\section{Outreach and media}
\cventry{2026}{Media engagement}{}{}{}{Interviewed by CTV News, with related coverage in \textit{The Guardian}, on black hole accretion and jet formation.}
\cventry{2024}{Media engagement}{}{}{}{For our \href{https://ui.adsabs.harvard.edu/abs/2024MNRAS.531..366M/abstract}{\textcolor{blue}{paper}} titled ``Continuum emission from within the plunging region of black hole discs'' I undertook a round of media interviews (including two \href{https://www.youtube.com/watch?v=49kfebgbFpA}{\textcolor{blue}{Youtube}} \href{https://www.youtube.com/shorts/JZ9qjB2PVuE}{\textcolor{blue}{interviews}} gaining over half a million views), to amplify the impact of our work. This paper ended up being the 60th most-discussed paper in the journals history.  
In total over 200 pieces of coverage reaching a global audience in excess of 150 million people were published, including the following:  
\newline \href{https://www.scientificamerican.com/article/black-holes-mysterious-plunging-region-matches-einsteins-predictions/}{\textcolor{blue}{Scientific American}}, 
\href{https://www.newscientist.com/article/2431520-einstein-was-right-about-the-way-matter-plunges-into-black-holes/}{\textcolor{blue}{New Scientist}}, 
\href{https://www.thetimes.co.uk/article/observed-at-last-the-black-hole-plunge-that-settles-einstein-v-newton-b208ql02s}{\textcolor{blue}{The Times}}, 
\href{https://edition.cnn.com/2024/05/17/world/black-holes-einstein-plunging-region-scn/index.html}{\textcolor{blue}{CNN}}, 
\href{https://www.dailymail.co.uk/sciencetech/article-13424999/Scientists-prove-Einstein-black- hole-theory.html}{\textcolor{blue}{The Daily Mail}}, 
\href{https://timesofindia.indiatimes.com/home/science/einstein-proved-right-study-gives-first-proof-of-plunging-regions-around-black-holes/articleshow/110185013.cms}{\textcolor{blue}{The Times of India}}, 
\href{https://nasa.einnews.com/pr_news/711733074/first-proof-that-plunging-regions-exist-around-black-holes-in-space}{\textcolor{blue}{NASA Today}}, 
\href{https://www.msn.com/en-us/news/technology/einstein-was-right-about-the-way-matter-plunges-into-black-holes/ar-BB1msVcN?ocid=BingNewsSearch}{\textcolor{blue}{MSN News}},
\href{https://www.independent.co.uk/space/scientists-albert-einstein-oxford-earth-b2545880.html}{\textcolor{blue}{The Independent}}, 
\href{https://www.msn.com/en-us/news/other/scientists-prove-that-plunging-regions-exist-around-black-holes-in-space/ar-BB1msCyz?ocid=BingNewsSearch}{\textcolor{blue}{The Evening Standard}}, 
\href{https://www.miragenews.com/oxford-confirms-existence-of-plunging-regions-1237487/}{\textcolor{blue}{Mirage News}}, 
\href{https://apnews.com/press-release/ein-presswire-newsmatics/oxford-7e62e22790c5de24d7d6f198b646f2f5}{\textcolor{blue}{Associated Press}}, 
\href{https://www.skyatnightmagazine.com/news/black-hole-plunge-region}{\textcolor{blue}{BBC}}
and 
\href{https://www.newsweek.com/black-holes-plunging-region-einstein-prediction-proven-correct-1901510}{\textcolor{blue}{Newsweek}}. 
}
\cventry{2023--2024}{Public lectures}{}{}{}{
\textit{-- “The spaghettification of stars by supermassive black holes: understanding one of nature’s most extreme events”} Public lecture for the event series “A Saturday Morning of Theoretical Physics” \href{https://podcasts.ox.ac.uk/spaghettification-stars-supermassive-black-holes-understanding-one-natures-most-extreme-events}{\textcolor{blue}{(podcast available)}}.
\newline 
\textit{-- ``Taking the plunge: detecting light from the very edge of black holes''} Public lecture to the Guildford Astronomy Society.}
\cventry{2018--2022}{Outreach -- UNIQ summer school tutor}{}{}{}{I ran problem classes for underprivileged high school students from around the UK who visit Oxford for a week, with the aim of encouraging them to apply to Universities in the future}{}

\section{Software}
\cvitem{FitTeD}{\href{https://github.com/fittingtransientswithdiscs}{\textcolor{blue}{GitHub organisation}} for fitting transients with relativistic discs. Personal repositories are being migrated here.}
\cvitem{FitTeD package}{Python package for light curves and spectra of evolving relativistic accretion discs, including \textbf{TDEFLARE} for rapid black hole masses from TDE optical/UV flares. \href{https://github.com/fittingtransientswithdiscs/FitTeD}{\textcolor{blue}{code}}, \href{https://arxiv.org/abs/2408.15048}{\textcolor{blue}{paper}}.}
\cvitem{tidalspin}{Bayesian black hole spin constraints from the Hills-mass limit in TDEs. \href{https://github.com/andymummeryastro/tidalspin}{\textcolor{blue}{code}}, \href{https://arxiv.org/abs/2312.00557}{\textcolor{blue}{paper}}.}
\cvitem{XSPEC models}{Public spectral models: \href{https://github.com/andymummeryastro/fullkerr}{\textcolor{blue}{fullkerr}}, \href{https://github.com/andymummeryastro/TDEdiscXraySpectrum}{\textcolor{blue}{TDEdiscXraySpectrum}}, \href{https://github.com/andymummeryastro/superkerr}{\textcolor{blue}{superkerr}}, \href{https://github.com/andymummeryastro/superkerrline}{\textcolor{blue}{superkerrline}}, \href{https://github.com/andymummeryastro/GR_prop_fluc}{\textcolor{blue}{GR\_prop\_fluc}}.}


\section{Publications}
\cvitem{}{76 publications; 42 first-author publications. Papers with A.~Goodwin or J.~Chakraborty marked $^{*}$ are joint first-author. $\dagger$~denotes a paper led by a student. Full, current lists: \href{https://ui.adsabs.harvard.edu/user/libraries/2zyx78XmSL6WL1D6o946hA}{\textcolor{blue}{NASA ADS}}, \href{https://scholar.google.com/citations?user=vEMN2XIAAAAJ\&hl=en}{\textcolor{blue}{Google Scholar}}, and \href{https://arxiv.org/a/mummery_a_1}{\textcolor{blue}{arXiv}}. Listed in reverse chronological order.}

\cvitem{[76].}{Chakraborty$^{*}$, \textbf{Mummery}$^{*}$, Quataert, Arcodia, Linial \& Kara (2026), Quasi-periodic eruption spectral-timing: EMRIs crossing warped accretion disks. Submitted. \href{https://arxiv.org/abs/2610.10110}{\textcolor{blue}{arXiv}}. $^{*}$Joint first authors.}
\cvitem{[75].}{$\dagger$Yu, \textbf{Mummery} \& Guolo (2026), HLX--1 as a repeating partial tidal disruption event around an intermediate-mass black hole. Accepted in ApJ. \href{https://arxiv.org/abs/2609.38694}{\textcolor{blue}{arXiv}}}
\cvitem{[74].}{\textbf{Mummery} \& Quataert (2026), Observational signatures of warped accretion flows in tidal disruption events. Submitted to ApJ. \href{https://arxiv.org/abs/2609.29787}{\textcolor{blue}{arXiv}}}
\cvitem{[73].}{\textbf{Mummery} \& Quataert (2026), The chirping of Lense--Thirring precession in tidal disruption event accretion flows. Submitted to ApJ. \href{https://arxiv.org/abs/2609.22559}{\textcolor{blue}{arXiv}}}
\cvitem{[72].}{$\dagger$Rule, \textbf{Mummery}, Balbus, Stone \& Zhang (2026), The plunging region of thin accretion discs across the black hole spin range. Accepted in MNRAS. \href{https://arxiv.org/abs/2608.06278}{\textcolor{blue}{arXiv}}}
\cvitem{[71].}{Goodwin, Miller-Jones, Sfaradi, \textbf{Mummery} et al. (2026), One year of broadband radio monitoring of the enigmatic transient GRB 250702B reveals the evolution of the relativistic jet. Submitted to MNRAS. \href{https://arxiv.org/abs/2608.03205}{\textcolor{blue}{arXiv}}}
\cvitem{[70].}{\textbf{Mummery}, Goodwin, Christy, Alexander \& Franz (2026), Triple radio flares from tidal disruption events: jet--wind collisions and the discovery of a third radio flare from AT2020vwl. Submitted. \href{https://arxiv.org/abs/2607.15438}{\textcolor{blue}{arXiv}}}
\cvitem{[69].}{Li, Christy, Alexander, Sfaradi, Huang, Jiang, Laskar, \textbf{Mummery} et al. (2026), VLA observations confirm AT 2023mfm as an off-nuclear tidal disruption event. RNAAS 10, 153. \href{https://arxiv.org/abs/2606.06595}{\textcolor{blue}{arXiv}}}
\cvitem{[68].}{Arcodia, Miniutti, Chakraborty, Franchini, Giustini, Linial, \textbf{Mummery} et al. (2026), Even a precessing clock is right twice per orbit --- the super-periods of eRO-QPE2 and challenges for QPE orbital models. ApJ 1003, 148. \href{https://arxiv.org/abs/2604.09788}{\textcolor{blue}{arXiv}}}
\cvitem{[67].}{Goodwin$^{*}$ \& \textbf{Mummery}$^{*}$ (2026), Resolving the sub-parsec circumnuclear density profiles of quiescent galaxies: evidence for Bondi accretion flows in TDE hosts. Submitted to ApJ. \href{https://arxiv.org/abs/2603.05412}{\textcolor{blue}{arXiv}}. $^{*}$Joint first authors.}
\cvitem{[66].}{Yao, Chornock, \textbf{Mummery} et al. (2026), AT2024lhc and AT2024kmq in the landscape of featureless tidal disruption events. MNRAS 549, stag920. \href{https://arxiv.org/abs/2602.21624}{\textcolor{blue}{arXiv}}}
\cvitem{[65].}{\textbf{Mummery} (2026), Accretion disks in (repeating) partial tidal disruption events: rapid state transitions, UV plateaus and flares from disk--remnant collisions. MNRAS 547, stag387. \href{https://arxiv.org/abs/2602.18698}{\textcolor{blue}{arXiv}}}
\cvitem{[64].}{Chakraborty, Masterson, \textbf{Mummery}, Kara, Panagiotou, Arcodia \& Berger (2026), X-ray spectral-timing properties of tidal disruption events. ApJ 1000, 95. \href{https://arxiv.org/abs/2602.16868}{\textcolor{blue}{arXiv}}}
\cvitem{[63].}{Goodwin$^{*}$ \& \textbf{Mummery}$^{*}$ (2026), A universal critical accretion rate for black hole jet formation. Nature Astronomy. \href{https://arxiv.org/abs/2602.14838}{\textcolor{blue}{arXiv}}. $^{*}$Joint first authors.}
\cvitem{[62].}{$\dagger$Guolo, \textbf{Mummery}, van Velzen et al. (2026), The wandering supermassive black hole powering the off-nuclear TDE AT2024tvd. ApJL 1004, L25. \href{https://arxiv.org/abs/2602.12272}{\textcolor{blue}{arXiv}}}
\cvitem{[61].}{Hamilton, \textbf{Mummery} \& Bland-Hawthorn (2026), On the radial velocity wave in the Galactic disk. Submitted to ApJ. \href{https://arxiv.org/abs/2602.06182}{\textcolor{blue}{arXiv}}}
\cvitem{[60].}{\textbf{Mummery} (2026), Supermassive black hole mass inference with the optical flares of tidal disruption events --- TDEFLARE. Submitted. \href{https://arxiv.org/abs/2601.14483}{\textcolor{blue}{arXiv}}}
\cvitem{[59].}{\textbf{Mummery} (2026), Black hole accretion disks with outflows II. Time-dependent Green's function solutions in Newtonian gravity. Submitted. \href{https://arxiv.org/abs/2601.11754}{\textcolor{blue}{arXiv}}}
\cvitem{[58].}{Omand, Sarin, Lamb, Perley, \textbf{Mummery} et al. (2026), Multiwavelength modeling of the luminous fast blue optical transient AT2024wpp. MNRAS 550, stag1104. \href{https://arxiv.org/abs/2601.03372}{\textcolor{blue}{arXiv}}}
\cvitem{[57].}{$\dagger$Xin \& \textbf{Mummery} (2026), The relativistic tidal tensor: general solutions for stationary axisymmetric spacetimes and the Hills mass of naked singularities. Physical Review D 113, 044036. \href{https://arxiv.org/abs/2511.21499}{\textcolor{blue}{arXiv}}}
\cvitem{[56].}{Chrimes, Jonker, Levan \& \textbf{Mummery} (2026), Luminous fast blue optical transients as very massive star core-collapse events. A\&A 706, A327. \href{https://arxiv.org/abs/2510.03402}{\textcolor{blue}{arXiv}}}
\cvitem{[55].}{\textbf{Mummery}, Metzger, van Velzen \& Guolo (2025), Tidal disruption event calorimetry: observational constraints on the physics of TDE optical flares. Submitted. \href{https://arxiv.org/abs/2512.09143}{\textcolor{blue}{arXiv}}}
\cvitem{[54].}{$\dagger$Guolo, \textbf{Mummery}, van Velzen et al. (2025), Compact accretion disks in the aftermath of tidal disruption events: parameter inference from joint X-ray spectra and UV/optical photometry. Submitted to ApJ. \href{https://arxiv.org/abs/2510.26774}{\textcolor{blue}{arXiv}}}
\cvitem{[53].}{Inkenhaag, Levan, \textbf{Mummery} \& Jonker (2025), AT 2018cow at $\sim$5 years: additional evidence for a tidal disruption origin. MNRAS 544, L108. \href{https://arxiv.org/abs/2510.08505}{\textcolor{blue}{arXiv}}}
\cvitem{[52].}{Saxton, Wevers, van Velzen, Alexander, Liu, \textbf{Mummery} et al. (2025), Rapid onset of a Comptonisation zone in the repeating tidal disruption event XMMSL2 J140446.9-251135. A\&A 704, A165. \href{https://arxiv.org/abs/2510.02905}{\textcolor{blue}{arXiv}}}
\cvitem{[51].}{Ramsden, Nicholl, McGee \& \textbf{Mummery} (2025), Evidence for a steeper SMBH--bulge mass relationship extended to low masses using TDE host galaxies. MNRAS 541, 1218. \href{https://arxiv.org/abs/2506.16155}{\textcolor{blue}{arXiv}}}
\cvitem{[50].}{\textbf{Mummery} \& Silk (2025), Black hole supercolliders. Physical Review Letters 134, 221401. \href{https://arxiv.org/abs/2506.01437}{\textcolor{blue}{arXiv}}}
\cvitem{[49].}{$\dagger$Rule, \textbf{Mummery}, Balbus, Stone \& Zhang (2025), The plunging region of a thin accretion disc around a Schwarzschild black hole. MNRAS 542, 377. \href{https://arxiv.org/abs/2505.13701}{\textcolor{blue}{arXiv}}}
\cvitem{[48].}{\textbf{Mummery}, Jiang, Ingram, Fabian \& Rule (2025), A rapid black hole spin or emission from the plunging region? MNRAS 544, 2880. \href{https://arxiv.org/abs/2505.13119}{\textcolor{blue}{arXiv}}}
\cvitem{[47].}{\textbf{Mummery} (2025), The X-ray variability of tidal disruption events. MNRAS 544, 1627. \href{https://arxiv.org/abs/2505.09238}{\textcolor{blue}{arXiv}}}
\cvitem{[46].}{\textbf{Mummery} (2025), Collisions with tidal disruption event disks: implications for quasi-periodic X-ray eruptions. Under review, MNRAS. \href{https://arxiv.org/abs/2504.21456}{\textcolor{blue}{arXiv}}}
\cvitem{[45].}{$\dagger$Guolo, \textbf{Mummery}, Ingram, Nicholl, Gezari \& Nathan (2025), A time-dependent solution for GSN 069 disk evolution and the nature of long-lived tidal disruption events. ApJ 992, 114. \href{https://arxiv.org/abs/2504.20148}{\textcolor{blue}{arXiv}}}
\cvitem{[44].}{Chakraborty, Kara, Arcodia, Buchner, Giustini, Hernández-García, Linial, Masterson, Miniutti, \textbf{Mummery} et al. (2025), Discovery of quasi-periodic eruptions in the TDE and extreme coronal line emitter AT2022upj. ApJL 983, L39. \href{https://arxiv.org/abs/2503.19013}{\textcolor{blue}{arXiv}}}
\cvitem{[43].}{\textbf{Mummery}, Guolo, Matthews, Newsome, Lintott \& Keel (2025), Galaxy-scale consequences of tidal disruption events: extended emission line regions, extreme coronal lines and infrared-to-optical light echoes. MNRAS 544, 2262. \href{https://arxiv.org/abs/2503.14163}{\textcolor{blue}{arXiv}}}
\cvitem{[42].}{\textbf{Mummery} (2025), Black hole--disc coevolution in the presence of magnetic fields: refining the Thorne limit with emission from within the plunging region. MNRAS 537, 1963. \href{https://arxiv.org/abs/2501.05843}{\textcolor{blue}{arXiv}}}
\cvitem{[41].}{Wevers, Guolo, Lockwood, \textbf{Mummery}, Pasham \& Arcodia (2025), Time-resolved Hubble Space Telescope UV observations of an X-ray quasi-periodic eruption source. ApJL 980, L1. \href{https://arxiv.org/abs/2501.03335}{\textcolor{blue}{arXiv}}}
\cvitem{[40].}{$\dagger$Guolo, \textbf{Mummery}, Wevers, Nicholl, Gezari, Ingram \& Pasham (2025), The properties of the GSN 069 accretion disk from a joint X-ray and UV spectral analysis: stress-testing quasi-periodic eruption models. ApJ 985, 146. \href{https://arxiv.org/abs/2501.03333}{\textcolor{blue}{arXiv}}}
\cvitem{[39].}{Witte \& \textbf{Mummery} (2025), Stepping up superradiance constraints on axions. Physical Review D 111, 083044. \href{https://arxiv.org/abs/2412.03655}{\textcolor{blue}{arXiv}}}
\cvitem{[38].}{Goodwin, \textbf{Mummery}, Laskar, Alexander et al. (2025), A second radio flare from the tidal disruption event AT2020vwl: a delayed outflow ejection? ApJ 981, 122. \href{https://arxiv.org/abs/2410.18665}{\textcolor{blue}{arXiv}}}
\cvitem{[37].}{\textbf{Mummery} \& van Velzen (2025), The optical, UV-plateau and X-ray tidal disruption event luminosity functions reproduced from first principles. MNRAS 541, 429. \href{https://arxiv.org/abs/2410.17087}{\textcolor{blue}{arXiv}}}
\cvitem{[36].}{Bright, Carotenuto, Fender, Choza, \textbf{Mummery} et al. (2025), The radio counterpart to the fast X-ray transient EP240414a. ApJ 981, 48. \href{https://arxiv.org/abs/2409.19055}{\textcolor{blue}{arXiv}}}
\cvitem{[35].}{$\dagger$Guolo \& \textbf{Mummery} (2025), The size of accretion disks from self-consistent X-ray spectra and UV/optical/NIR photometry fitting: ASASSN-14li and HLX-1 --- diskSED. ApJ 978, 167. \href{https://arxiv.org/abs/2408.17296}{\textcolor{blue}{arXiv}}}
\cvitem{[34].}{\textbf{Mummery}, Nathan, Ingram \& Gardner (2025), Fitting transients with discs --- FitTeD. MNRAS 544, 2225. \href{https://arxiv.org/abs/2408.15048}{\textcolor{blue}{arXiv}}}
\cvitem{[33].}{Nicholl, Pasham, \textbf{Mummery} et al. (2024), Quasi-periodic X-ray eruptions years after a nearby tidal disruption event. Nature 634, 804. \href{https://arxiv.org/abs/2409.02181}{\textcolor{blue}{arXiv}}}
\cvitem{[32].}{\textbf{Mummery} \& Stone (2024), The three-dimensional structure of black hole accretion flows within the plunging region. MNRAS 532, 3395. \href{https://arxiv.org/abs/2407.02164}{\textcolor{blue}{arXiv}}}
\cvitem{[31].}{\textbf{Mummery}, Jiang \& Fabian (2024), Plunging region emission in the X-ray binary MAXI J0637-430. MNRAS 533, L83. \href{https://arxiv.org/abs/2406.14957}{\textcolor{blue}{arXiv}}}
\cvitem{[30].}{Chrimes, Coppejans, Jonker, Levan, Groot, \textbf{Mummery} \& Stanway (2024), Multi-wavelength observations of the luminous fast blue optical transient AT2023fhn. A\&A 691, A329. \href{https://arxiv.org/abs/2406.13821}{\textcolor{blue}{arXiv}}}
\cvitem{[29].}{\textbf{Mummery}, Ingram, Davis \& Fabian (2024), Continuum emission from within the plunging region of black hole discs. MNRAS 531, 366. \href{https://arxiv.org/abs/2405.09175}{\textcolor{blue}{arXiv}}}
\cvitem{[28].}{\textbf{Mummery} \& Turner (2024), The turbulent variability of accretion discs observed at high energies. MNRAS 530, 4730. \href{https://arxiv.org/abs/2404.09564}{\textcolor{blue}{arXiv}}}
\cvitem{[27].}{Pasham, Tombesi, Suková, Zajaček, Rakshit, Coughlin, Kosec, Karas, Masterson, \textbf{Mummery} et al. (2024), A case for a binary black hole system revealed via quasi-periodic outflows. Science Advances 10, eadj8898. \href{https://arxiv.org/abs/2402.10140}{\textcolor{blue}{arXiv}}}
\cvitem{[26].}{\textbf{Mummery}, Mori \& Balbus (2024), The dynamics of accretion flows near to the innermost stable circular orbit. MNRAS 529, 1900. \href{https://arxiv.org/abs/2402.05561}{\textcolor{blue}{arXiv}}}
\cvitem{[25].}{\textbf{Mummery} \& Ingram (2024), Reflecting on naked singularities: iron line fitting as a probe of the cosmic censorship conjecture. MNRAS 528, 2015. \href{https://arxiv.org/abs/2401.10565}{\textcolor{blue}{arXiv}}}
\cvitem{[24].}{\textbf{Mummery} (2024), The maximum mass of a black hole which can tidally disrupt a star: measuring black hole spins with tidal disruption events --- tidalspin. MNRAS 527, 6233. \href{https://arxiv.org/abs/2312.00557}{\textcolor{blue}{arXiv}}}
\cvitem{[23].}{\textbf{Mummery}, Balbus \& Ingram (2024), Testing theories of accretion and gravity with super-extremal Kerr discs. MNRAS 527, 5956. \href{https://arxiv.org/abs/2311.15742}{\textcolor{blue}{arXiv}}}
\cvitem{[22].}{\textbf{Mummery}, van Velzen, Nathan, Ingram, Hammerstein, Fraser-Taliente \& Balbus (2024), Fundamental scaling relationships revealed in the optical light curves of tidal disruption events. MNRAS 527, 2452. \href{https://arxiv.org/abs/2308.08255}{\textcolor{blue}{arXiv}}}
\cvitem{[21].}{Chrimes, Jonker, Levan, Coppejans, Gaspari, Gompertz, Groot, Malesani, \textbf{Mummery} et al. (2024), AT2023fhn (the Finch): a luminous fast blue optical transient at a large offset from its host galaxy. MNRAS 527, L47. \href{https://arxiv.org/abs/2307.01771}{\textcolor{blue}{arXiv}}}
\cvitem{[20].}{Inkenhaag, Jonker, Levan, Chrimes, \textbf{Mummery}, Perley \& Tanvir (2023), Late-time HST UV and optical observations of AT 2018cow: extracting a cow from its background. MNRAS 525, 4042. \href{https://arxiv.org/abs/2308.07381}{\textcolor{blue}{arXiv}}}
\cvitem{[19].}{\textbf{Mummery} (2023), Extending the theory of propagating fluctuations: the first fully relativistic treatment and analytical Fourier--Green's functions. MNRAS 523, 3629. \href{https://arxiv.org/abs/2305.18314}{\textcolor{blue}{arXiv}}}
\cvitem{[18].}{Wiseman, Wang, Hönig, Castro-Segura, Clark, Frohmaier, Fulton, Leloudas, Middleton, Müller-Bravo, \textbf{Mummery} et al. (2023), Multiwavelength observations of the extraordinary accretion event AT2021lwx. MNRAS 522, 3992. \href{https://arxiv.org/abs/2303.04412}{\textcolor{blue}{arXiv}}}
\cvitem{[17].}{\textbf{Mummery} \& Balbus (2023), Accretion within the innermost stable circular orbit: analytical thermodynamic solutions in the adiabatic limit. MNRAS 521, 2439. \href{https://arxiv.org/abs/2302.14437}{\textcolor{blue}{arXiv}}}
\cvitem{[16].}{\textbf{Mummery} \& Balbus (2023), A complete characterisation of the orbital shapes of the non-circular Kerr geodesic solutions with circular orbit constants of motion. Physical Review D 107, 124058. \href{https://arxiv.org/abs/2302.01159}{\textcolor{blue}{arXiv}}}
\cvitem{[15].}{\textbf{Mummery}, Wevers, Saxton \& Pasham (2023), From X-rays to physical parameters: a comprehensive analysis of thermal tidal disruption event X-ray spectra. MNRAS 519, 5828. \href{https://arxiv.org/abs/2301.07419}{\textcolor{blue}{arXiv}}}
\cvitem{[14].}{\textbf{Mummery} (2023), Asymptotic Green's function solutions of the general relativistic thin disc equations. MNRAS 518, 1905. \href{https://arxiv.org/abs/2210.06161}{\textcolor{blue}{arXiv}}}
\cvitem{[13].}{\textbf{Mummery} \& Balbus (2022), The high energy probability distribution of accretion disc luminosity fluctuations. MNRAS 517, 3423. \href{https://arxiv.org/abs/2210.01450}{\textcolor{blue}{arXiv}}}
\cvitem{[12].}{\textbf{Mummery} \& Balbus (2022), Inspirals from the innermost stable circular orbit of Kerr black holes: exact solutions and universal radial flow. Physical Review Letters 129, 161101. \href{https://arxiv.org/abs/2209.03579}{\textcolor{blue}{arXiv}}}
\cvitem{[11].}{Goodwin, van Velzen, Miller-Jones, \textbf{Mummery} et al. (2022), AT2019azh: an unusually long-lived, radio-bright thermal tidal disruption event. MNRAS 511, 5328. \href{https://arxiv.org/abs/2201.03744}{\textcolor{blue}{arXiv}}}
\cvitem{[10].}{\textbf{Mummery} (2021), Tidal disruption event discs are larger than they seem: removing systematic biases in TDE X-ray spectral modelling. MNRAS 507, L24. \href{https://arxiv.org/abs/2108.10160}{\textcolor{blue}{arXiv}}}
\cvitem{[9].}{\textbf{Mummery} (2021), A unified model of tidal destruction events in the disc-dominated phase. arXiv. \href{https://arxiv.org/abs/2104.06212}{\textcolor{blue}{arXiv}}}
\cvitem{[8].}{\textbf{Mummery} (2021), A maximum X-ray luminosity scale of disc-dominated tidal disruption events. MNRAS 504, 5144. \href{https://arxiv.org/abs/2104.06203}{\textcolor{blue}{arXiv}}}
\cvitem{[7].}{\textbf{Mummery} \& Balbus (2021), Hard X-ray emission from a Compton scattering corona in large black hole mass tidal disruption events. MNRAS 504, 4730. \href{https://arxiv.org/abs/2104.06195}{\textcolor{blue}{arXiv}}}
\cvitem{[6].}{\textbf{Mummery} \& Balbus (2021), An upper observable black hole mass scale for tidal disruption events with thermal X-ray spectra. MNRAS 505, 1629. \href{https://arxiv.org/abs/2104.06177}{\textcolor{blue}{arXiv}}}
\cvitem{[5].}{\textbf{Mummery} \& Balbus (2020), ASASSN-15lh: a TDE about a maximally rotating $10^9 M_\odot$ black hole. MNRAS 497, L13. \href{https://arxiv.org/abs/2006.00803}{\textcolor{blue}{arXiv}}}
\cvitem{[4].}{\textbf{Mummery} \& Balbus (2020), The spectral evolution of disc-dominated tidal disruption events. MNRAS 492, 5655. \href{https://arxiv.org/abs/1912.06577}{\textcolor{blue}{arXiv}}}
\cvitem{[3].}{\textbf{Mummery} \& Balbus (2019), Evolution of relativistic thin discs with a finite ISCO stress: II. Late time behaviour. MNRAS 489, 143. \href{https://arxiv.org/abs/1908.00348}{\textcolor{blue}{arXiv}}}
\cvitem{[2].}{\textbf{Mummery} \& Balbus (2019), Evolution of relativistic thin discs with a finite ISCO stress: I. Stalled accretion. MNRAS 489, 132. \href{https://arxiv.org/abs/1908.00322}{\textcolor{blue}{arXiv}}}
\cvitem{[1].}{Balbus \& \textbf{Mummery} (2018), The evolution of Kerr discs and late-time tidal disruption event light curves. MNRAS 481, 3348. \href{https://arxiv.org/abs/1809.02146}{\textcolor{blue}{arXiv}}}

\section{Selected invited talks (2026)}
\cvitem{October}{Astronomy colloquium, University of California, Berkeley, USA.}
\cvitem{August}{Chandra Dynamic Universe meeting, Boston, USA.}
\cvitem{August}{COSPAR, Florence, Italy.}
\cvitem{August}{TDE conference, Crete, Greece.}
\cvitem{July}{Dutch Black Hole Consortium meeting, Amsterdam, the Netherlands.}
\cvitem{June}{X-ray Universe meeting, Elche, Spain.}
\cvitem{April}{Perimeter Institute, Waterloo, Canada.}
\cvitem{April}{New York University, New York, USA.}
\cvitem{February}{TDE workshop, Kyoto, Japan.}

% Older talks kept for a future highlights rewrite.
\iffalse
\section{Talks}
% To add – Edingburough, Radboud, KITP plasma, Microquasar, SW17, QPE Madrid – invitations. 

\cvitem{$\bullet$}{ \textit{“The three-dimensional structure of black hole accretion flows within the plunging region”}
\newline Plasma physics talk series — Oxford University 
\newline High energy astrophysics talk series — Radboud University  }

\cvitem{$\bullet$}{ \textit{“Continuum emission from within the plunging region of black hole discs”}
\newline High energy astrophysics talk series — Radboud University  
\newline ICAP — Institute of Astronomy Paris  
\newline Making waves and instabilities (conference) — University of Oxford } 

\cvitem{$\bullet$}{\textit{“The maximum mass of a black hole which can tidally disrupt a star"}\newline
Dynamics talk series --- Institute for Advanced Study}

\cvitem{$\bullet$}{ \textit{“Super-extremal Kerr discs as a probe of cosmic censorship”}
\newline BlackStar talk series — Caltech } 

\cvitem{$\bullet$}{\textit{“Theoretical models of tidal disruption events”}\newline
European astronomical society annual meeting — Krakow \textbf{(Invited review talk)}
\newline Tidal disruption event conference — Radboud University \textbf{(Invited review talk)}}

\cvitem{$\bullet$}{\textit{“Tidal disruption events as a time domain astronomy $M-\sigma$ relationship”}\newline
Origins \textbf{Colloquium} -- University of Virginia  \newline
Astrophysics \textbf{Colloquium} -- Liverpool John Moores  \newline
Tidal Disruption Events: entering the data rich era --- Heraklion \newline
Anticipating the Rising Tide of Tidal Disruption Events --- KITP UCSB \newline
High energy astrophysics talk series --- Flatiron institue \newline 
Bahcall lunch --- Institute for Advanced Study 
\newline 
Journal club — Newcastle University 
\newline 
High energy talk series — Birmingham University 
\newline 
Astrophysics \textbf{Colloquium} — Oxford University   
\newline
Galaxies series  — Leiden University                
\newline
Tidal disruption event conference — Radboud University   
}

\cvitem{$\bullet$}{\textit{“Accretion within the ISCO of Kerr black holes”}\newline
High energy astrophysics talk series  — Leiden University                                
\newline 
European astronomical society annual meeting — Krakow 
\newline 
Astrophysics \textbf{Colloquium}  — Radboud University   \newline 
Astrophysical fluids seminar series  — Cambridge University (DAMTP)  
\newline 
High energy talk series — Birmingham University 
}                            

\cvitem{$\bullet$}{\textit{“A relativistic theory of propagating fluctuations: implications for black hole variability”}\newline
European astronomical society annual meeting — Krakow \newline 
High energy astrophysics talk series  — Amsterdam University}

\cvitem{$\bullet$}{\textit{“Inspirals onto Kerr black holes: new exact solutions of Einsteins equations”}\newline
Theoretical Physics \textbf{Colloquium}   — Oxford University}

\cvitem{$\bullet$}{\textit{“A unified model of tidal disruption events”}\newline 
X-ray astrophysics talk series — Southampton University 
\newline
High energy astrophysics talk series  — Princeton University                                               
\newline 
National Astronomy Meeting  — University of Bath}
 
 \cvitem{$\bullet$}{\textit{“A relativistic disc model of the tidal disruption event ASASSN-14li”}\newline
Summer talk series  — University of Cambridge  	 
\newline
Compact Objects talk series — Computational Centre for Astrophysics (New York)
\newline
Astrophysical Transients talk series — New York University 
\newline
Tidal disruption events: confronting theory with observations — Yukawa Institute, Kyoto
\newline
TEXAS symposium on relativistic astrophysics — University of Portsmouth                    	   
\newline
SPIMAX talk series — University of Oxford	
}

 \cvitem{$\bullet$}{\textit{“A time-dependent, general relativistic, theory of thin accretion discs” }\newline
SPIMAX talk series — University of Oxford \newline
Planets, Stars and Discs Conference — University of Oxford }
\fi

% References omitted from the public website CV.
\iffalse
\section{References}

\begin{tabular}{lr}
% Referee 1
\begin{minipage}[t]{3in}
\textbf{Prof. Steven Balbus}\\
\textit{Savilian Chair of Astrophysics, \\ Department of Physics,} \\
University of Oxford\\
\Letter\ \href{mailto:steven.balbus@physics.ox.ac.uk}{steven.balbus@physics.ox.ac.uk}
\end{minipage}
&
% Referee 2
\begin{minipage}[t]{3in}
\textbf{Prof. Chris Reynolds}\\
\textit{Professor of Astronomy,} \\
\textit{Director,  Joint Space Science Institute, }\\
University of Maryland, College Park\\
\Letter\ \href{mailto:creynold@umd.edu}{creynold@umd.edu}
\end{minipage}
\\
\\ % Additional newline for spacing.
% Referee 3
\begin{minipage}[t]{3in}
\textbf{Prof. James Stone}\\
\textit{Professor of Astrophysics, \\
School of Natural Sciences}\\
{Institute for Advanced Study} \\
\Letter\ \href{mailto:jmstone@ias.edu}{jmstone@ias.edu}
\end{minipage}
&
% Referee 4
\begin{minipage}[t]{3in}
\textbf{Prof. Rob Fender}\\
\textit{Professor of Astronomy,} \\
\textit{Department of Physics}\\
University of Oxford\\
\Letter\ \href{mailto:rob.fender@physics.ox.ac.uk}{rob.fender@physics.ox.ac.uk}
\end{minipage}
\\
\\ % Additional newline for spacing.
% Referee 3
\begin{minipage}[t]{3in}
\textbf{Prof. Andrew Fabian}\\
\textit{Professor of Astrophysics, \\
Institute for Astronomy} \\
{University of Cambridge}\\
\Letter\ \href{mailto:acf@ast.cam.ac.uk}{acf@ast.cam.ac.uk}
\end{minipage}
&
% Referee 4
\begin{minipage}[t]{3in}
\textbf{Prof. Sjoert van Velzen}\\
\textit{Assistant Professor of Astronomy,} \\
\textit{Department of Physics}\\
Leiden University\\
\Letter\ \href{mailto:sjoert@strw.leidenuniv.nl}{sjoert@strw.leidenuniv.nl}
\end{minipage}
\\
\end{tabular}
\fi

\end{document}